\section*{Limitations}
\label{sec:limitations}

\paragraph{Minimal two-owner setting.}
Our benchmark restricts collaboration to scenarios involving exactly two owners and two corresponding agents.
This minimal configuration enables controlled analysis and clear attribution of constraint violations, but it does not capture additional coordination challenges that may arise in larger groups.
In particular, multi-owner settings may introduce emergent dynamics such as coalition formation, indirect information leakage, or shifting responsibility boundaries, which are beyond the scope of this work.

\paragraph{Natural language constraints and automated evaluation.}
Privacy constraints in our benchmark are expressed in natural language and evaluated using an LLM-based judge.
While this design allows scalable and domain-flexible assessment, it may be sensitive to ambiguity in constraint interpretation.
More structured representations of privacy constraints, as well as hybrid evaluation schemes combining automated judgment with targeted human verification, remain promising directions for future work.

\paragraph{Fixed privacy constraints.}
Our benchmark assumes fixed, owner-defined privacy constraints throughout each collaboration episode. This design choice enables controlled and systematic measurement of how privacy restrictions affect collaborative reasoning and coordination. In particular, by keeping constraints fixed, PAC-BENCH isolates the interaction between privacy constraints and multi-agent collaboration in an analyzable manner. However, real-world privacy preferences may evolve during interaction, and agents may need to support dynamic privacy negotiation or incremental disclosure. Extending the benchmark to incorporate such dynamic or negotiated privacy policies is a natural next step and a promising direction for future research on private-agent collaboration.

\paragraph{Fixed agent prompting and memory formulation.}
Agent behavior in our experiments is conditioned on a fixed prompting strategy and memory formulation.
As a result, some observed failure modes may reflect limitations of current prompting approaches rather than fundamental limits of constraint-aware collaboration.
Exploring alternative agent architectures, such as explicit constraint reasoning modules or stricter memory separation mechanisms, is left for future investigation.

\paragraph{Collaboration scope.}
Our evaluation focuses on 20-turn collaborative tasks with a predefined goal.
Longer-horizon interactions, where agents must consistently enforce owner-defined constraints across extended interaction histories, may exhibit qualitatively different behaviors.
Studying such long-term collaborations is necessary to fully characterize the challenges of ownership-aware agent systems in real-world deployments.


\section*{Ethics Statements}

\paragraph{Use of data and privacy protection.}
This work studies multi-agent collaboration under explicit privacy constraints, with the goal of evaluating how agents balance task completion and privacy preservation. To minimize ethical risks, the benchmark and experiments are carefully designed to avoid the use of real personal or sensitive data.

Although the benchmark focuses on privacy-aware collaboration, no real-world personal data is used. All agent memories, scenarios, and interaction contexts are synthetically generated. The privacy constraints imposed on agents are not derived from actual user records, but are instead constructed based on publicly available standards, prior academic literature, and documented domain practices, which are properly cited in the paper. As a result, the experimental setup does not involve personally identifiable information (PII), private communications, or proprietary datasets.

\paragraph{Construction of privacy constraints.}
The privacy constraints in \ours are designed to model realistic restrictions that may arise in practical deployments, such as limits on information disclosure or communication scope. Importantly, these constraints are abstracted representations grounded in established confidentiality and privacy frameworks, rather than reflections of specific individuals or organizations.

This abstraction is a deliberate design choice that enables systematic study of privacy-constrained collaboration while minimizing the risk of privacy leakage or re-identification.

\paragraph{Human evaluation.}
We conduct a limited human evaluation to validate the reliability of the automated privacy compliance metrics. The authors serve as human annotators, assessing agent messages solely based on observable outputs and predefined evaluation rubrics. Annotators do not have access to any sensitive information, and the evaluation task does not expose them to harmful, personal, or distressing content.

\paragraph{Scope and ethical trade-offs.}
While PAC-BENCH captures key aspects of privacy-aware collaboration, the use of synthetic data and abstracted constraints may limit direct applicability to real-world deployments involving actual user data. We view this as an intentional ethical trade-off that prioritizes safety, reproducibility, and controlled analysis while enabling principled evaluation of privacy-constrained interaction dynamics.

Overall, this work aims to support responsible research on collaborative agent systems by providing a benchmark that facilitates the study of privacy-aware behavior without relying on real personal data or introducing avoidable ethical risks.
